% Generated draft: requires mathematical and source review.
\documentclass[11pt]{article}
\usepackage[margin=25mm]{geometry}
\usepackage{fontspec}
\setmainfont{Latin Modern Roman}
\usepackage{amsmath,amssymb,mathtools,unicode-math}
\setmathfont{Latin Modern Math}
\usepackage{enumitem,needspace,xurl,hyperref}
\setlength{\parindent}{0pt}
\setlength{\parskip}{5pt}
\setlength{\emergencystretch}{3em}
\newcommand{\N}{\mathbb{N}}
\newcommand{\Z}{\mathbb{Z}}
\newcommand{\Q}{\mathbb{Q}}
\newcommand{\R}{\mathbb{R}}
\newcommand{\C}{\mathbb{C}}
\newcommand{\F}{\mathbb{F}}
\newcommand{\A}{\mathbb{A}}
\newcommand{\PP}{\mathbb{P}}
\newcommand{\E}{\mathbb{E}}
\newcommand{\eps}{\varepsilon}
\newcommand{\dd}{\,\mathrm{d}}
\begin{document}
\section*{UnsolvedMath ID: 30005752}
\textbf{Source problem code:} catalogue entry
\subsection{Definitions and mathematical statement}
The supplied clean statement concerns a quantum greedy basis with parameters \(\delta_1,\delta_2\). The extracted report appears to display the factorization
\[
\Delta=\begin{pmatrix}\delta_1&0\\0&\delta_2\end{pmatrix},\qquad
\Omega=\begin{pmatrix}0&-1\\1&0\end{pmatrix},\qquad
B=\Delta\Omega=\begin{pmatrix}0&-\delta_1\\\delta_2&0\end{pmatrix}.
\]
This is a reconstruction of damaged matrix formatting and requires confirmation against the report. The record does not explicitly specify the domain of the \(\delta_i\); positive integers are a natural candidate domain for the divisibility condition, but this is an additional interpretation.

\textbf{Intended conjectural implication, with unresolved conventions.} If \(\mathcal G_q(B)\) denotes the quantum greedy basis intended by the supplied clean statement, the requested implication is
\[
\bigl(\delta_1\mid\delta_2\ \text{or}\ \delta_2\mid\delta_1\bigr)
\quad\Longrightarrow\quad
\mathcal G_q(B)\text{ has the intended positivity property}.
\]
The construction of \(\mathcal G_q(B)\), its coefficient ring, the compatible quantum data, and the meaning of positivity are not defined in the supplied excerpt. Consequently this implication is not yet a completely specified mathematical conjecture.

One possible interpretation is universal Laurent positivity: every basis element, when expanded in the quantum Laurent monomials of every cluster, has coefficients that are finite sums of allowed quantum-parameter monomials with nonnegative integer coefficients. Under that interpretation, the intended quantifiers would be: for every admissible pair \((\delta_1,\delta_2)\) satisfying the divisibility condition, every element of the corresponding quantum greedy basis is positive in every cluster. This interpretation, including the admissible pairs, coefficient conventions, and monomial normalization, requires source confirmation. The phrase ``positivity is preserved'' does not itself identify an operation under which positivity is preserved.

The original excerpt instead places the conjecture in a discussion of cluster scattering diagrams. It introduces \(N=\Z^2\), a set \(N^+\) of positive elements whose exact definition is not supplied, and the skew-symmetric form determined by \(\Omega\). It describes a Lie algebra with generators \(X_n\), \(n\in N^+\), and relations
\[
[X_n,X_{n'}]=[\{n,n'\}]_qX_{n+n'},
\qquad [a]_q=\frac{q^a-q^{-a}}{q-q^{-1}}.
\]
The displayed formula for \([a]_q\) is the apparent restoration of corrupted extraction. The excerpt then uses a degree completion \(\widehat{\mathfrak g}_q\), the BCH exponential group \(G_q=\exp(\widehat{\mathfrak g}_q)\), and, for positive rational parameters \(a,b\), elements apparently given by
\[
\Psi_{a,b}[n]=\exp\!\left(\sum_{j=1}^{\infty}
\frac{(-1)^{j+1}q^{jb}}{j[ja]_q}X_{jn}\right).
\]
The coefficient field, interpretation of rational powers of \(q\), degree, and completion conventions are not specified in the record. The excerpt states that if \(\{n_2,n_1\}=c>0\), then, for ``any'' \(b_i\),
\[
\Psi_{c,b_2}[n_2]\Psi_{c,b_1}[n_1]
=
\Psi_{c,b_1}[n_1]\Psi_{c,b_1+b_2}[n_1+n_2]\Psi_{c,b_2}[n_2].
\]
The scope of ``any'' is unclear relative to the earlier requirement that \(b\) be positive rational. The excerpt says that the classical case is obtained as \(q\to1\), with \(\Psi_{a,b}[n]\) reducing to \(\Psi[n]^{1/a}\), and that the same divisibility condition emerges through the pentagon relation. It does not explicitly identify the positivity in this discussion with positivity of a quantum greedy basis; that connection remains to be checked.
\subsection{Short English statement}
The intended question asks whether the quantum greedy basis is positive whenever either parameter divides the other. A precise final formulation requires confirmation of the parameter domain, quantum cluster algebra conventions, and meaning of positivity. The original excerpt discusses scattering diagrams without explicitly identifying its positivity property with quantum greedy-basis positivity.
\subsection{Sources}
\begin{itemize}
\item Supplied local UnsolvedMath registry record, permanent ID \(30005752\), title \emph{Positivity Criterion for Quantum Greedy Bases}, exact source problem code \texttt{catalogue entry}; Oberwolfach Reports Open Problems collection, proposed year 2024. No proposer is named in the record.
\item \emph{Cluster Algebras and Its Applications} (2024), Oberwolfach Reports, pp.\ 69--136, DOI \texttt{10.4171/owr/2024/2}. Supplied source report and literature reference ``original OWR report'': \url{https://doi.org/10.4171/owr/2024/2}. The extracted discussion carries the talk title \emph{Classical vs Quantum Pentagon Relations} and the name Tomoki Nakanishi; this does not establish authorship of the conjecture, which is attributed there to reference [3]. The scattering-diagram method is associated with reference [1]. Neither numbered reference is resolved by the supplied excerpt.
\item Supplied literature reference ``quantum greedy-basis paper'': \url{https://doi.org/10.1016/j.aim.2016.03.022}.
\item Supplied literature reference ``quantum scattering-diagram paper'': \url{https://arxiv.org/abs/2202.01588}.
\item The exact supplied literature-check date is \texttt{2026-08-22}. The supplied dated curation labels the problem partially solved and assesses that the quantum greedy basis exists, positivity fails in general, positive cases are known, and the stated divisibility criterion remains conjectural. This is supplied dated curation derived from the OWR collection, not an independently verified literature assessment. Its time-sensitive status and cited claims require independent verification.
\item The registry labels the statement \texttt{corrected\_verified} and reports that the official report was checked directly and extraction artifacts were removed. These are supplied provenance claims; no independent source visit or verification is asserted in this draft.
\end{itemize}
\end{document}
